\documentclass[10pt,a4paper]{amsart}
\usepackage[margin=19mm]{geometry}
\usepackage{amsmath,amssymb,mathtools}
\usepackage[scr=rsfs,cal=euler]{mathalfa}
\usepackage{xfrac}
\usepackage[unicode,hidelinks,bookmarks=false]{hyperref}
\newcommand{\ie}{i.e.\ }
\newcommand{\cf}{cf.\ }
\newcommand{\resp}{respectively\ }
\newcommand{\Subsection}{\subsection}
\providecommand{\nobreakdash}{\nobreak\hbox{-}}
\setlength{\emergencystretch}{2em}
\multlinegap0pt
\usepackage{amssymb}
\usepackage{algorithm}\usepackage{algpseudocode}
\usepackage{tikz}
\usepackage{xcolor}
\usepackage{mathtools}
\newtheorem{lemma}{Lemma}
\title{Locally Semi-Equitable Colourings of BIBDs}
\author{Andrea C. Burgess, William Kellough, David A. Pike}
\date{Source: arXiv:2605.20002v1 (CC BY 4.0)}
\begin{document}
\maketitle

\noindent\emph{Source and licence.} Locally Semi-Equitable Colourings of BIBDs. Version arXiv:2605.20002v1 (CC BY 4.0), distributed by arXiv under the Creative Commons Attribution 4.0 International licence, http://creativecommons.org/licenses/by/4.0/, as recorded in the arXiv licence metadata for this exact version and shown on the abstract page https://arxiv.org/abs/2605.20002v1. Full licence notice: \texttt{licenses/arxiv-2605.20002-CC-BY-4.0.txt}.\par\smallskip

\noindent\textit{Editorial scope.} Counted unit: the article's lemma alpha (source0.tex lines 193--195). The article's complete original proof of it is delivered with it (source0.tex lines 197--231, one \texttt{proof} environment of 35 lines). No mathematics is written, repaired or re-derived: the statement and the proof are the article's own lines, in the article's own notation, and no retained line is substituted or re-flowed.  No further source-local context is needed: every cross-reference of every retained line resolves inside the two delivered spans.  Nothing was omitted from any retained span, so the package carries no omission record.  No retained line contains a citation, so the wrapper carries no bibliography and no bibliography entry.  No retained line contains a figure environment, an image-inclusion command or drawing code, so no image is delivered and the package declares no asset dependency.  Editorial formatting only, in addition: an AMS \texttt{amsart} preamble that restates the article's own declarations and macros used by the retained lines, namely the environments \texttt{lemma} on separate counters, together with the standard packages those lines need.  The retained environments print in this document as Lemma 1 in order of appearance, whereas the article prints the same items with the numbering of its own full document.  The complete native source of the article as distributed in the arXiv e-print for this exact version is bundled unchanged as \texttt{sources/200960/source0.tex}.
\begin{lemma}\label{lem: l-1 divides k, colour classes the same size, calculating alpha}
    In a $(v,k,\lambda)$-BIBD with a $0$-ULSE $\ell$-colouring, each colour appears in exactly \[\frac{r^2}{\lambda} - \frac{(\ell - 1) r^2}{\lambda k} + \frac{(\ell - 1)r}{k}\] blocks and each colour class has size $\frac{v}{\ell}$.
\end{lemma}

\begin{proof}

Let $q = \frac{k}{\ell - 1}$. Let $c_i$ be a colour with corresponding colour class $C_i$.

Since each colour appears either $q$ or $0$ times in a block,

\begin{equation}\label{eq: k=0 mod l-1, alpha + gamma = b}
    \alpha_i + \gamma_i = b.
\end{equation}

By counting, in two different ways, the number of pairs $(x,B)$ where $x$ is a point coloured $c_i$ and $B$ is a block containing $x$ we obtain the following.

\begin{equation}\label{eq: k=0 mod l-1, alpha q = rCi}
    \alpha_i q = r|C_i|
\end{equation}

By counting, in two different ways, the number of pairs $(xy, B)$ where $xy$ is an unordered pair of points both coloured $c_i$ and $B$ is a block containing the pair $xy$ we obtain the following.

\begin{equation}\label{eq: k=0 mod l-1, alpha qC2 = lambda |C_i|C2}
    \alpha_i \binom{q}{2} = \lambda \binom{|C_i|}{2}
\end{equation}

By combining Equations \eqref{eq: k=0 mod l-1, alpha q = rCi} and \eqref{eq: k=0 mod l-1, alpha qC2 = lambda |C_i|C2} we obtain \[r|C_i| \frac{q-1}{2} = \frac{\lambda}{2}|C_i| (|C_i| - 1).\] Solving for $|C_i|$ gives 
\begin{equation}\label{eq: formula for |Ci| in terms of krlambda}
    |C_i| = \frac{kr}{\lambda (\ell -1)} - \frac{r}{\lambda} + 1.
\end{equation}

Since the right-hand side is independent of $i$, every colour class has the same size. Therefore, $|C_i| = \frac{v}{\ell}$ for every colour class $C_i$. Substituting $|C_i| = \frac{kr}{\lambda (\ell -1)} - \frac{r}{\lambda} + 1$ into Equations \eqref{eq: k=0 mod l-1, alpha + gamma = b} and \eqref{eq: k=0 mod l-1, alpha q = rCi} gives the following values for $\alpha_i$ and $\gamma_i$.

\begin{align}
    \alpha_i &= \frac{r^2}{\lambda} - \frac{(\ell - 1) r^2}{\lambda k} + \frac{(\ell - 1)r}{k} \label{eq: alpha formula} \\
    \gamma_i &= b - \frac{r^2}{\lambda} + \frac{(\ell - 1) r^2}{\lambda k} - \frac{(\ell - 1)r}{k} \label{eq: gamma formula}
\end{align}

\end{proof}

\end{document}
